\documentclass[a4paper,11pt]{article}
\usepackage[margin=1in]{geometry}
\usepackage[T1]{fontenc}
\usepackage[utf8]{inputenc}
\usepackage{lmodern}
\usepackage{amsmath,amssymb}
\usepackage{booktabs}
\usepackage{graphicx}
\usepackage{array}
\usepackage{hyperref}

\renewcommand{\thefigure}{S\arabic{figure}}
\renewcommand{\thetable}{S\arabic{table}}
\setcounter{figure}{0}
\setcounter{table}{0}

\begin{document}
\begin{center}
{\LARGE\bfseries Online Resource 1}\par
\vspace{0.4em}
{\large\bfseries Scientometrics}\par
\vspace{0.8em}
{\large Linking Undergraduate Innovation and Entrepreneurship Training Projects to Academic-Unit Publication Portfolios:\\
A Cross-Lingual Measure of Research--Training Alignment}\par
\vspace{0.8em}
[Author Name]\\
School of Remote Sensing and Information Engineering, Wuhan University, Wuhan, China\\
Corresponding author: [email to be completed]
\end{center}
\vspace{1em}

\section*{Supplementary Figures}


% ---------------------------------------------------------
% Figure S1
% ---------------------------------------------------------

\begin{figure}[htbp]
\centering
\includegraphics[width=0.82\linewidth]{FigS1.pdf}
\caption{Observed supervisory-load distribution among 1,834 individual advisors after multi-advisor project fields were split using the same rules as the main analysis. Counts are based on distinct advisor--project links in the 2020--2024 registry. The 2--3-project category is the complement of the one-project and $\ge4$-project groups; mean load is 2.314 projects, median 2, and maximum 10}
\label{fig:s1}
\end{figure}


% ---------------------------------------------------------
% Figure S2
% ---------------------------------------------------------

\begin{figure}[htbp]
\centering
\includegraphics[width=0.98\linewidth]{FigS2.pdf}
\caption{Four-quadrant summary for clustered documents. Bars count documents assigned to the 886 non-outlier topics; BERTopic outliers are excluded from the quadrant counts. The prevalence thresholds used to classify topics are defined against the full article and project corpora, as specified in the main text}
\label{fig:s2}
\end{figure}


% ---------------------------------------------------------
% Figure S3
% ---------------------------------------------------------

\begin{figure}[htbp]
\centering
\includegraphics[width=0.92\linewidth]{FigS3.pdf}
\caption{Topic dynamics in UIETP projects, 2020--2024: aggregate share of the top 15 clustered topics by project-side prevalence versus all other clustered topics. Annual denominators include all projects, whereas BERTopic outliers are absent from both plotted numerators; consequently, the two lines do not sum to 100\%. This is a prevalence-share comparison, not a formal diversity index}\label{fig:s3}
\end{figure}


% ---------------------------------------------------------
% Figure S4
% ---------------------------------------------------------

\begin{figure}[htbp]
\centering
\includegraphics[width=0.92\linewidth]{FigS4.pdf}
\caption{Top-15 project-side topic prevalence by year, 2020--2024}\label{fig:s4}
\end{figure}


% ---------------------------------------------------------
% Figure S5
% ---------------------------------------------------------

\begin{figure}[htbp]
\centering
\includegraphics[height=0.80\textheight,keepaspectratio]{FigS5.pdf}
\caption{Descriptive RTAS means for all 40 academic units using the full project dataset. The one-way academic-unit ANOVA shown in the figure uses 39 units because the single-project academic unit cannot contribute a within-group variance estimate; this ANOVA effect size is distinct from the adjusted mixed-model ICC in the main text. Darker blue indicates higher descriptive RTAS and lighter blue lower RTAS. The ordering is descriptive and should not be interpreted as a performance ranking}
\label{fig:s5}
\end{figure}


\clearpage
\section*{Supplementary Tables}

% ---------------------------------------------------------
% Table S1
% ---------------------------------------------------------

The official 2021--2024 approval lists classify projects as innovation training or entrepreneurship training/practice. Matching on (year, normalized title) identifies 158 unique entrepreneurship projects inside the frozen 3,714-row dataset (4.3\%; 74/23/26/35 across 2021--2024; 53 university-, 64 provincial-, and 41 national-level). These projects have lower alignment than the remaining 3,556 projects: mean RTAS $=0.100$ (SD $=0.053$) versus $0.135$ (SD $=0.073$); Welch $t=-7.90$, $p<0.001$, Cohen $d=-0.49$ (pooled SD). Table~S1 re-estimates the frozen random-intercept specification (REML, clustered by academic unit, covariates taken from the frozen dataset unchanged) after excluding these 158 projects: the complete-case sample falls from 3,231 to 3,085 projects, ICC changes from 0.7376 to 0.7373, every coefficient retains its sign and significance classification, and the administrative-tier indicators remain non-significant. The 2020 cohort is retained in both fits because no type-labeled official registry exists for that year.

\begin{table}[htbp]
\centering
\caption{Sensitivity of the primary mixed model (the primary mixed-model table in the main manuscript) to excluding the 158 officially typed entrepreneurship projects (2021--2024)}
\label{tab:s1}
\small
\begin{tabular}{p{0.36\linewidth}rrrr}
\toprule
& \multicolumn{2}{c}{Primary ($N=3{,}231$)} & \multicolumn{2}{c}{Excl. entrepreneurship ($N=3{,}085$)}\\
\cmidrule(lr){2-3}\cmidrule(lr){4-5}
Predictor & $\beta$ & $p$ & $\beta$ & $p$\\
\midrule
Provincial vs university & 0.00062 & 0.720 & 0.00085 & 0.627\\
National vs university & $-0.00162$ & 0.452 & $-0.00205$ & 0.357\\
Year (per year; centered 2022) & 0.00313 & $<0.001$ & 0.00293 & $<0.001$\\
Advisor prior 3y publications, $\log(1+x)$ & 0.00427 & $<0.001$ & 0.00480 & $<0.001$\\
Advisor prior supervised projects, $\log(1+x)$ & $-0.00629$ & $<0.001$ & $-0.00567$ & $<0.001$\\
\midrule
ICC & 0.7376 & & 0.7373 & \\
Number of academic units & 39 & & 39 & \\
\bottomrule
\end{tabular}

\vspace{1mm}
\footnotesize Same specification as the primary mixed-model table in the main manuscript: RTAS $\sim$ provincial $+$ national $+$ year $+$ $\log(1+\text{prior 3y publications})$ $+$ $\log(1+\text{prior supervised projects})$, random intercept by academic unit, REML. Entrepreneurship projects are identified from the official 2021--2024 project-type field; 2020 track membership is unobservable, so no 2020 project is excluded.
\end{table}

% ---------------------------------------------------------
% Table S2
% ---------------------------------------------------------

The archived official 2022 approval list contains no university-level entries (316 provincial and 150 national projects; 435 projects from that list remain in the frozen dataset after the eight-label study-frame exclusion). To check whether the RQ1 administrative-tier comparison or the primary mixed model depends on this structurally different cohort, Table~S2 repeats both analyses after dropping all 435 projects approved in 2022. Because 2022 contributes no university-level projects, the university-tier mean is necessarily unchanged. At RQ1 ($N=3{,}279$ projects with RTAS), the overall contrast strengthens slightly ($F(2,3276)=35.80$, $p<0.001$, $\eta^2=2.14\%$); both significant Tukey contrasts against university tier remain significant, and the national--provincial contrast remains non-significant ($+0.0053$, adjusted $p=0.31$). In the mixed model, complete cases fall from 3,231 to 2,852 and the ICC is 0.7409; every coefficient retains its sign and significance classification, and the administrative-tier indicators remain non-significant.

\begin{table}[htbp]
\centering
\caption{Sensitivity to excluding the entire 2022 cohort, whose archived official list contains no university-level records}
\label{tab:s2}
\footnotesize
\begin{tabular}{lcc}
\toprule
\multicolumn{3}{l}{\textit{Panel A. RQ1 project-level RTAS by administrative tier}}\\
& Primary & Excluding 2022\\
& ($N=3{,}714$) & ($N=3{,}279$)\\
\midrule
University, mean (SD) & 0.1212 (0.0694) & 0.1212 (0.0694)\\
Provincial, mean (SD) & 0.1391 (0.0728) & 0.1407 (0.0728)\\
National, mean (SD) & 0.1460 (0.0731) & 0.1460 (0.0723)\\
One-way ANOVA $F$ & 35.72 & 35.80\\
ANOVA $\eta^2$ & 1.89\% & 2.14\%\\
Tukey: provincial $-$ university & $+0.0179^{***}$ & $+0.0195^{***}$\\
Tukey: national $-$ university & $+0.0248^{***}$ & $+0.0248^{***}$\\
Tukey: national $-$ provincial & $+0.0070$ ($p=0.082$) & $+0.0053$ ($p=0.312$)\\
\bottomrule
\end{tabular}

\vspace{1.5em}

\begin{tabular}{lrrrr}
\toprule
\multicolumn{5}{l}{\textit{Panel B. Primary mixed model (random intercept by academic unit, REML)}}\\
& \multicolumn{2}{c}{Primary ($N=3{,}231$)} & \multicolumn{2}{c}{Excluding 2022 ($N=2{,}852$)}\\
\cmidrule(lr){2-3}\cmidrule(lr){4-5}
Predictor & $\beta$ & $p$ & $\beta$ & $p$\\
\midrule
Provincial vs university & 0.00062 & 0.720 & 0.00166 & 0.351\\
National vs university & $-0.00162$ & 0.452 & $-0.00212$ & 0.353\\
Year (per year; centered 2022) & 0.00313 & $<0.001$ & 0.00285 & $<0.001$\\
Advisor prior 3y publications, $\log(1+x)$ & 0.00427 & $<0.001$ & 0.00434 & $<0.001$\\
Advisor prior supervised projects, $\log(1+x)$ & $-0.00629$ & $<0.001$ & $-0.00494$ & 0.002\\
\midrule
ICC & 0.7376 & & 0.7409 & \\
Number of academic units & 39 & & 39 & \\
\bottomrule
\end{tabular}

\vspace{1mm}
\footnotesize $^{***}p<0.001$. Panel~A uses project-level RTAS values (one-way ANOVA, Tukey HSD); Panel~B uses the same specification as the primary mixed-model table in the main manuscript and Table~S1. The 2022 cohort is the only year with no university-level records in the archived source registry; excluding it cannot alter the university-tier mean but changes the composition of the other two tiers.
\end{table}

\end{document}
